% =====================================================================
%  Fibonacci Character Sums at Rank Index Two and Jacobsthal Sums
%  Author: Majid Ghandali  (ORCID: 0009-0001-1097-1770)
%  Source of truth: ghandali_index2.tex + ghandali_index2.bib
%  See README.md for build conventions
% =====================================================================
\documentclass[11pt]{amsart}
\usepackage[margin=1in]{geometry}
\usepackage{amsmath,amssymb,amsthm,mathtools}
\usepackage[dvipsnames]{xcolor}
\usepackage{booktabs}
\usepackage{enumitem}
\usepackage[numbers,sort&compress]{natbib}
\usepackage[colorlinks=true,linkcolor=blue,citecolor=blue,urlcolor=blue]{hyperref}

% ---------- theorem environments ----------
\theoremstyle{plain}
\newtheorem{theorem}{Theorem}[section]
\newtheorem{proposition}[theorem]{Proposition}
\newtheorem{lemma}[theorem]{Lemma}
\newtheorem{corollary}[theorem]{Corollary}
\newtheorem{conjecture}[theorem]{Conjecture}
\newtheorem{mainthm}{Theorem}
\renewcommand{\themainthm}{\Alph{mainthm}}
\theoremstyle{definition}
\newtheorem{definition}[theorem]{Definition}
\newtheorem{example}[theorem]{Example}
\newtheorem{remark}[theorem]{Remark}
\newtheorem{problem}[theorem]{Problem}
\newtheorem{observation}[theorem]{Observation}

% ---------- macros ----------
\newcommand{\Fp}{\mathbb{F}_p}
\newcommand{\Fpp}{\mathbb{F}_{p^2}}
\newcommand{\leg}[2]{\left(\frac{#1}{#2}\right)}
\newcommand{\qleg}[2]{\left(\frac{#1}{#2}\right)_{\!4}}
\newcommand{\Jac}{\mathcal{J}}
% Note: previously also defined as \phi_2; kept only the calligraphic form used throughout the text.

\begin{document}

\title[Index-two Fibonacci character sums]{Fibonacci Character Sums at Rank Index Two\\and Jacobsthal Sums}

\author{Majid Ghandali}
\address{Independent Researcher, Tehran, Iran}
\email{majid.ghandali@gmail.com}
\thanks{ORCID: 0009-0001-1097-1770.}
\date{}
\subjclass[2020]{11B39, 11A15, 11L10}
\keywords{Fibonacci numbers, rank of apparition, Legendre symbol, character sums, Jacobsthal sums, sums of two squares, quartic residues}

\begin{abstract}
Let $p>5$ be prime, let $z(p)$ be the rank of apparition of $p$
in the Fibonacci sequence, and put
\[
  S(p)=\sum_{j=1}^{z(p)}\leg{F_j}{p},
  \qquad
  i_p=\frac{p-\leg{5}{p}}{z(p)}.
\]
Earlier work of the author evaluates $S(p)$ in the full-rank case
$i_p=1$.  This paper treats the next layer, namely the rank-index-two
case $i_p=2$.

We first show that $i_p=2$ forces $p\equiv1\pmod4$.  In the
$p\equiv1\pmod8$ branches, a change of variables through a primitive
root in the split case and a Cayley parametrisation of the norm-one
torus in the inert case reduce $S(p)$ to classical Jacobsthal sums; the
complementary $p\equiv5\pmod8$ branches are handled by reflection and
character identities.

Writing
\[
  \Jac(c)=\sum_{x\in\mathbb F_p}\leg{x}{p}\leg{x^2+c}{p},
\]
we obtain the following explicit laws.  In the split case, with
\[
  \delta_p=
  \begin{cases}
  1,&p\equiv1\pmod5,\\
  -1,&p\equiv4\pmod5,
  \end{cases}
\]
and with $p=a^2+b^2$, $a$ odd, $b>0$ even, and
$a^{*}\in\{\pm a\}$ determined by $a^{*}\equiv1\pmod4$, one has
\[
  S(p)=
  \begin{cases}
    \delta_p a^{*},&p\equiv1\pmod8,\\
    -\delta_p,&p\equiv5\pmod8.
  \end{cases}
\]
In the inert case one has
\[
  S(p)=
  \begin{cases}
    \displaystyle \frac12\Jac(5)\in\{\pm b\},&p\equiv1\pmod8,\\[4pt]
    0,&p\equiv5\pmod8,
  \end{cases}
\]
and in the first branch the sign is determined by the congruence
$S(p)\equiv -a^{*}\pmod5$.

A companion verifier checks the theorem-level identities and records
the numerical data used in the paper; the source code and
machine-readable outputs are available in the versioned reproducibility
archive~\cite{ghandaliIndex2Archive}.
\end{abstract}

\maketitle

% =====================================================================
\section{Introduction}\label{sec:intro}
% =====================================================================

\subsection{Setting}
For a prime $p>5$ let $\chi=\chi_p=\leg{\cdot}{p}$, $\varepsilon_p=\leg5p$,
$n_p=p-\varepsilon_p$, let $z(p)$ be the rank of apparition of $p$ in $(F_n)$, and
\[
  i_p:=\frac{n_p}{z(p)}\in\mathbb Z_{\ge1},\qquad
  S(p):=\sum_{j=1}^{z(p)}\chi(F_j).
\]
Earlier work \cite{Ghandali2026Preprint} --- the \emph{full-rank paper} --- gives an
explicit evaluation of $S(p)$ for $i_p=1$; for split primes it reads
$S(p)=\delta_p$, where $\delta_p\in\{\pm1\}$ and
$\delta_p\equiv p\pmod 5$. The density of the primes with $i_p=1$ is,
under GRH, a special case of Chen's theorem on primitive points of
one-dimensional tori \cite{Chen2002}: the total density is Artin's
constant $A$, with inert and split branches $10A/19$ and $9A/19$
respectively (see the conditional corollary in \cite{Ghandali2026Preprint}). The
maximal-order problem for units in real quadratic fields was initiated
numerically by Ishikawa--Kitaoka \cite{IshikawaKitaoka1998} and settled
conditionally by Chen--Kitaoka--Yu \cite{ChenKitaokaYu2000} and Roskam
\cite{Roskam2000}; for the fundamental unit $\varphi$ (rather than the
norm-one element $\eta=-\varphi^2$), Roskam's inert value is
$20A/19$. See Lenstra \cite{Lenstra1977} and Murty \cite{Murty1983}
for the general conditional framework. The present note asks what
happens one level down, at $i_p=2$.

\subsection{Main results}
Write $p=a^2+b^2$ with $a$ odd, $b$ even and $b>0$ whenever $p\equiv1\pmod4$, and let
$a^{*}\in\{\pm a\}$ be determined by $a^{*}\equiv1\pmod4$. Let
$\Jac(c)=\sum_{x\in\Fp}\chi(x)\chi(x^2+c)$ be the Jacobsthal sum.

\begin{mainthm}[Split primes, index two]\label{thm:A}
Let $p$ be split ($\varepsilon_p=1$) with $i_p=2$. Then $p\equiv1\pmod4$, and
\[
  S(p)=\begin{cases}\delta_p\,a^{*}, & p\equiv1\pmod8,\\[2pt] -\delta_p, & p\equiv5\pmod8.\end{cases}
\]
\end{mainthm}

\begin{mainthm}[Inert primes, index two]\label{thm:B}
Let $p$ be inert, that is $\varepsilon_p=\leg5p=-1$, and suppose
$i_p=2$. Then $p\equiv1\pmod4$. If $p\equiv5\pmod8$, then
\[
  S(p)=0.
\]
If $p\equiv1\pmod8$, write
\[
  p=(a^{*})^2+b^2,\qquad a^{*}\equiv1\pmod4,\qquad b>0.
\]
Then
\[
  S(p)=\frac12\Jac(5)=\frac12\Jac(-5),
\]
and the sign is given by
\[
  S(p)=
  \begin{cases}
    -b,& b\equiv a^{*}\pmod5,\\
    \phantom{-}b,& b\equiv -a^{*}\pmod5.
  \end{cases}
\]
Equivalently, when $p\equiv1\pmod8$, $S(p)$ is the unique element of
$\{\pm b\}$ congruent to $-a^{*}$ modulo~$5$.
\end{mainthm}

\begin{proposition}[Reflection vanishing]\label{prop:reflection-vanishing}
If $z(p)$ is odd, then $p\equiv1\pmod4$. Moreover, if in addition $p\equiv5\pmod8$,
then $S(p)=0$.
\end{proposition}

\subsection{What is classical and what is claimed}
\begin{center}
\begin{tabular}{@{}p{0.42\textwidth}p{0.32\textwidth}p{0.20\textwidth}@{}}
\toprule
Ingredient & Source & Role here\\
\midrule
$i_p$ even $\iff p\equiv1\pmod4$ & classical (Lehmer \cite{DHLehmer}) & Lemma~\ref{lem:index-parity}\\
Jacobsthal: $(\tfrac12\Jac(r))^2+(\tfrac12\Jac(n))^2=p$ & Jacobsthal \cite{Jacobsthal}, \cite{BEW} & input\\
$\Jac(-1)=\sum\chi(x^3-x)=-2a^{*}$ ($p\equiv1\bmod 8$) & Gauss; \cite[Ch.~6,9]{BEW} & input\\
$\qleg5p\,\chi(\varphi)=\delta_p$ & elementary lemma; proof supplied here & Lemma~\ref{lem:quartic-unit}\\
Rank multiplier / reflection & classical (Lucas \cite{Lucas1878}, Somer \cite{Somer1980}); see also \cite{Ghandali2026HalfPeriod} & Lemmas~\ref{lem:multiplier}, \ref{lem:reflection}\\
Half-Jacobsthal reduction ($p\equiv1\pmod8$, $i_p=2$) & this paper & Props.~\ref{prop:split-red}, \ref{prop:inert-red}\\
Sign of $\tfrac12\Jac(5)$ via biquadratic reciprocity & classical reciprocity + application here & Lemma~\ref{lem:sign5}, Thm.~\ref{thm:B}\\
Inert index-two Dirichlet density $=10A/19$ (GRH) & Roskam \cite{Roskam2000} + Chen \cite{Chen2002} + 2-adic dictionary (\S\ref{sec:context}) & Cor.~\ref{cor:inert-index-two-density}\\
Theorems~\ref{thm:A}, \ref{thm:B} & this paper & main results\\
Proposition~\ref{prop:reflection-vanishing} & classical; also used in \cite{Ghandali2026HalfPeriod} & structural\\
\bottomrule
\end{tabular}
\end{center}
\noindent
The contribution claimed here is the reduction of the Fibonacci
character sum at rank index two to the specific Jacobsthal sums
$\Jac(-1)$ and $\Jac(\pm5)$, together with the resulting explicit
formulas for $S(p)$.  The classical evaluation theory of Jacobsthal
sums and the relevant reciprocity laws are used as inputs; no new
general evaluation theory for Jacobsthal sums is claimed.
We have not located a prior reference recording these particular
rank-index-two Fibonacci reductions, nor their combination with the
quartic-reciprocity sign determination in the inert branch.  This is a
qualified literature statement, not a proof of priority.  The inert Dirichlet
density $10A/19$ recorded in Corollary~\ref{cor:inert-index-two-density}
below is derived from the GRH-conditional density theorems of Roskam and
Chen after translating their statements into the present rank-index
notation; it is included only as conditional context
(Remark~\ref{rem:natural-vs-dirichlet}).

\subsection{Organisation}
Section~\ref{sec:prelim} fixes notation. Section~\ref{sec:struct}
proves the structural lemmas. Sections~\ref{sec:split}
and~\ref{sec:inert} prove Theorems~\ref{thm:A} and~\ref{thm:B}.
Section~\ref{sec:numerics} records the numerical verification of the
theorem-level identities (including the signed inert law).
Section~\ref{sec:context} records a GRH-conditional density corollary
obtained from Roskam's and Chen's density theorems after translating
their maximal-order statements into the present rank-index notation.
This conditional discussion is included only for context and is not used
in the proofs of Theorems~\ref{thm:A} and~\ref{thm:B}.
Section~\ref{sec:open} records further problems and directions.

% =====================================================================
\section{Notation and preliminaries}\label{sec:prelim}
% =====================================================================

Let $\varphi=(1+\sqrt5)/2$, $\psi=(1-\sqrt5)/2$, so $\varphi\psi=-1$ and
$F_n=(\varphi^n-\psi^n)/(\varphi-\psi)$. Reducing modulo $p$ we regard
$\varphi,\psi\in\Fp$ if $\varepsilon_p=1$ and $\varphi,\psi\in\Fpp\setminus\Fp$ with
$\varphi^p=\psi$ if $\varepsilon_p=-1$. Put
\[
  \eta:=\varphi/\psi=-\varphi^2 .
\]
Then $F_n\equiv0\pmod p\iff\eta^n=1$, so $z(p)=\operatorname{ord}(\eta)$ and
$\eta^{n_p}=1$.

\begin{definition}
For $p\equiv1\pmod4$ fix $\iota\in\Fp$ with $\iota^2=-1$; then
$\chi(\iota)=(-1)^{(p-1)/4}=\leg2p$. For split $p\equiv1\pmod4$ put
$\qleg5p:=\chi(\varphi-\psi)\equiv5^{(p-1)/4}\pmod p$; this does not depend on the
choice of square root since $\chi(-1)=1$. For $p\equiv1\pmod8$ write $2=s^2$
in $\Fp$ and set $\qleg2p:=\chi(s)$; this is the quartic symbol of $2$.
\end{definition}

\begin{definition}[Jacobsthal sums]\label{def:jacobsthal-sums}
$\Jac(c):=\sum_{x\in\Fp}\chi(x)\chi(x^2+c)$ for $c\in\Fp^{\times}$.
\end{definition}

\begin{remark}
Throughout this paper, the term ``Jacobsthal sum'' refers to the
classical quadratic character sum
\[
  \Jac(c)=\sum_{x\in\Fp}\chi(x)\chi(x^2+c).
\]
It is unrelated to the Jacobsthal and Jacobsthal--Lucas recurrence
sequences commonly denoted by $J_n$ and $j_n$.
\end{remark}

\begin{lemma}[Jacobsthal]\label{lem:jacobsthal}
Let $p\equiv1\pmod4$, let $r\in\mathbb F_p^\times$ be a quadratic residue and
let $n\in\mathbb F_p^\times$ be a quadratic non-residue. Then
\[
  \frac12\Jac(r) \text{ is odd},\qquad
  \frac12\Jac(n) \text{ is even},
\]
and
\[
  \left(\frac12\Jac(r)\right)^2+
  \left(\frac12\Jac(n)\right)^2=p.
\]
Moreover, for every $c,u\in\mathbb F_p^\times$,
\[
  \Jac(cu^2)=\chi(u)\Jac(c).
\]
\end{lemma}

\begin{proof}
The first assertions are the classical theorem of Jacobsthal on the two
components in the representation of primes $p\equiv1\pmod4$ as a sum of two
squares; see Jacobsthal \cite{Jacobsthal} or
Berndt--Evans--Williams \cite[Chapter~6]{BEW}.

For the transformation formula, substitute $x=uy$:
\[
\begin{aligned}
 \Jac(cu^2)
 &=\sum_{x\in\mathbb F_p}\chi(x)\chi(x^2+cu^2)\\
 &=\sum_{y\in\mathbb F_p}\chi(uy)\chi(u^2(y^2+c))\\
 &=\chi(u)\sum_{y\in\mathbb F_p}\chi(y)\chi(y^2+c)\\
 &=\chi(u)\Jac(c).
\end{aligned}
\]
\end{proof}

\begin{lemma}[Jacobsthal congruence]\label{lem:jac-cong}
Let $p\equiv1\pmod4$, and write
\[
  p=a^2+b^2,\qquad a\ \text{odd},\quad b\ \text{even},
\]
with $a^{*}\in\{\pm a\}$ determined by $a^{*}\equiv1\pmod4$.
Let $c\in\mathbb F_p^\times$, and put $m=(p-1)/2$. Then
\[
  \Jac(c)\equiv -\binom{m}{m/2}\,c^{(p-1)/4}\pmod p.
\]
By the classical Gauss congruence
\[
  \binom{m}{m/2}\equiv 2a^{*}\pmod p,
\]
one obtains
\[
  \tfrac12\Jac(c)\equiv -a^{*}\,c^{(p-1)/4}\pmod p.
\]
In particular, for $p\equiv1\pmod8$,
\[
  \Jac(-1)=-2a^{*}.
\]
\end{lemma}

\begin{proof}
Working in $\mathbb F_p$, one has $\chi(x)\equiv x^m$ for all $x\in\mathbb F_p$,
with both sides equal to $0$ at $x=0$. Hence
\[
  \Jac(c)
  =
  \sum_{x\in\mathbb F_p}\chi(x)\chi(x^2+c)
  \equiv
  \sum_{x\in\mathbb F_p}x^m(x^2+c)^m
  \pmod p.
\]
Expanding the binomial gives
\[
  \sum_{k=0}^{m}\binom{m}{k}c^{m-k}
  \sum_{x\in\mathbb F_p}x^{m+2k}.
\]
Using the rule
\[
  \sum_{x\in\mathbb F_p}x^e\equiv
  \begin{cases}
    -1\pmod p,& (p-1)\mid e,\ e>0,\\
    0\pmod p,& \text{otherwise},
  \end{cases}
\]
only the middle term $k=m/2$ survives. Therefore
\[
  \Jac(c)\equiv
  -\binom{m}{m/2}c^{m/2}
  =
  -\binom{m}{m/2}c^{(p-1)/4}
  \pmod p.
\]
The evaluation of the central binomial coefficient is classical:
\[
  \binom{m}{m/2}\equiv 2a^{*}\pmod p;
\]
see, for example, Berndt--Evans--Williams~\cite[Chapters~6 and~9]{BEW}.
Thus
\[
  \tfrac12\Jac(c)\equiv -a^{*}\,c^{(p-1)/4}\pmod p.
\]

Now suppose $p\equiv1\pmod8$ and put $c=-1$. Then
\[
  (-1)^{(p-1)/4}=1,
\]
so
\[
  \tfrac12\Jac(-1)\equiv -a^{*}\pmod p.
\]
Since $-1$ is a quadratic residue, Lemma~\ref{lem:jacobsthal} implies that
$\tfrac12\Jac(-1)$ is the odd component in the representation
$p=a^2+b^2$. Hence
\[
  \tfrac12\Jac(-1)=\pm a^{*}.
\]
The congruence modulo $p$ determines the sign. Indeed, $2|a^{*}|<p$ for
$p>5$, so $a^{*}$ and $-a^{*}$ are not congruent modulo $p$. Therefore
\[
  \tfrac12\Jac(-1)=-a^{*},
\]
or equivalently
\[
  \Jac(-1)=-2a^{*}.
\]
\end{proof}

\begin{lemma}\label{lem:Jminus1}
Let $p\equiv1\pmod8$, and write
\[
  p=a^2+b^2,\qquad a\ \text{odd},\quad b\ \text{even},\quad b>0,
\]
with $a^{*}\in\{\pm a\}$ determined by $a^{*}\equiv1\pmod4$. Then
\[
  \Jac(-1)=\sum_{x\in\mathbb F_p}\chi(x^3-x)=-2a^{*}.
\]
\end{lemma}

\begin{proof}
This is the final assertion of Lemma~\ref{lem:jac-cong}.
\end{proof}

% =====================================================================
\section{Structural lemmas}\label{sec:struct}
% =====================================================================

\begin{lemma}[Parity]\label{lem:parity}
$S(p)\equiv z(p)-1\pmod 2$.
\end{lemma}
\begin{proof}
For $1\le j<z(p)$ one has $F_j\not\equiv0$, so $\chi(F_j)=\pm1$; and $\chi(F_{z(p)})=0$.
\end{proof}

\begin{lemma}[Index parity]\label{lem:index-parity}
$\eta^{n_p/2}=\chi(-1)$. Hence $i_p$ is even if and only if $p\equiv1\pmod4$.
\end{lemma}
\begin{proof}
If $\varepsilon_p=1$: $\eta^{(p-1)/2}=(-1)^{(p-1)/2}\varphi^{p-1}=(-1)^{(p-1)/2}$.
If $\varepsilon_p=-1$: $\varphi^{p+1}=\varphi\psi=-1$, so
$\eta^{(p+1)/2}=(-1)^{(p+1)/2}\varphi^{p+1}=(-1)^{(p-1)/2}$. Finally $i_p$ is even
iff $z(p)\mid n_p/2$ iff $\eta^{n_p/2}=1$.
\end{proof}

\begin{lemma}[Rank multiplier]\label{lem:multiplier}
Let $z=z(p)$. Then $F_{z+1}=\varphi^{z}\in\Fp$ and $F_{z+1}^2=(-1)^z$. Consequently
$\chi(F_{z+1})=1$ if $z$ is even, while if $z$ is odd then $p\equiv1\pmod4$ and
$\chi(F_{z+1})=\leg2p$.
\end{lemma}
\begin{proof}
$F_z=0$ means $\varphi^z=\psi^z=(-1)^z\varphi^{-z}$, so $\varphi^{2z}=(-1)^z$, and
$F_{z+1}=(\varphi^{z+1}-\psi^{z+1})/(\varphi-\psi)=\varphi^z$. If $z$ is odd,
$F_{z+1}^2=-1$ forces $p\equiv1\pmod4$ and $F_{z+1}=\pm\iota$.
\end{proof}

\begin{lemma}[Reflection]\label{lem:reflection}
For $0\le j\le z=z(p)$: $F_{z-j}=(-1)^{j+1}F_{z+1}F_j$.
\end{lemma}
\begin{proof}
Using $\psi^z=\varphi^z$ and $\varphi\psi=-1$:
$\varphi^{z-j}-\psi^{z-j}=\varphi^z(\varphi^{-j}-\psi^{-j})
=\varphi^z(-1)^{j}(\psi^{j}-\varphi^{j})$; divide by $\varphi-\psi$ and apply
Lemma~\ref{lem:multiplier}.
\end{proof}

\begin{proof}[Proof of Proposition~\ref{prop:reflection-vanishing}]
Let $z=z(p)$ be odd. Since $p$ is odd and
\[
  n_p=p-\varepsilon_p
\]
is even, the quotient
\[
  i_p=\frac{n_p}{z}
\]
is even. Hence Lemma~\ref{lem:index-parity} gives $p\equiv1\pmod4$, so
$\chi(-1)=1$. By Lemma~\ref{lem:reflection}, $j\mapsto z-j$ is a fixed-point-free
involution of $\{1,\dots,z-1\}$ with
$\chi(F_{z-j})=\chi(F_{z+1})\chi(F_j)$
(note that the factor $\chi((-1)^{j+1})$ is absorbed because $\chi(-1)=1$), and
$\chi(F_z)=0$. Hence $S(p)=\chi(F_{z+1})S(p)$, and Lemma~\ref{lem:multiplier} gives
$\chi(F_{z+1})=\leg2p$. In particular $S(p)=0$ whenever $\leg2p=-1$, i.e.,
$p\equiv5\pmod8$.
\end{proof}

\begin{corollary}\label{cor:vanishing}
$S(p)=0$ in each of the following cases: (a) $\varepsilon_p=-1$, $i_p=2$,
$p\equiv5\pmod8$; (b) $\varepsilon_p=1$, $i_p=4$, $p\equiv5\pmod8$.
\end{corollary}
\begin{proof}
In (a) $z=(p+1)/2$ and in (b) $z=(p-1)/4$; both are odd.
\end{proof}

\begin{lemma}\label{lem:split-index2-chi}
If $\varepsilon_p=1$ and $i_p=2$, then $\chi(\varphi)=-\leg2p$.
\end{lemma}
\begin{proof}
Here $p\equiv1\pmod4$ and $\operatorname{ord}\eta=(p-1)/2$ is even, so
$\eta^{(p-1)/4}=-1$. But
$\eta^{(p-1)/4}=(-1)^{(p-1)/4}\varphi^{(p-1)/2}=\leg2p\chi(\varphi)$.
\end{proof}

\begin{lemma}[Quartic character of the golden unit]\label{lem:quartic-unit}
Let $p>5$ be a split prime for $K=\mathbb Q(\sqrt5)$, and suppose
$p\equiv1\pmod4$. Define
\[
  \delta_p=
  \begin{cases}
    1,& p\equiv1\pmod5,\\
    -1,& p\equiv4\pmod5.
  \end{cases}
\]
Then
\[
  \qleg5p\,\chi(\varphi)=\delta_p.
\]
Equivalently,
\[
  \chi(\varphi+2)=\delta_p.
\]
\end{lemma}

\begin{proof}
Since
\[
  \varphi(\varphi-\psi)=\varphi\sqrt5=\varphi+2,
\]
we have
\[
  \qleg5p\,\chi(\varphi)
  =
  \chi(\sqrt5)\chi(\varphi)
  =
  \chi(\varphi+2).
\]
Thus it suffices to determine whether $\varphi+2$ is a square in $\mathbb F_p$.

The element $\varphi-1$ is a root of the quadratic
\[
  X^2+X-1=0.
\]
Let $\zeta\in\mathbb F_{p^2}$ be a primitive fifth root of unity.
Since $p\equiv\pm1\pmod5$, Frobenius acts by
\[
  \zeta\longmapsto\zeta^p\in\{\zeta,\zeta^{-1}\}.
\]
Hence $\zeta+\zeta^{-1}\in\mathbb F_p$, and the two elements
\[
  \zeta+\zeta^{-1},
  \qquad
  \zeta^2+\zeta^{-2}
\]
are precisely the roots $\varphi-1,\psi-1$ of $X^2+X-1$.
Replacing $\zeta$ by $\zeta^2$ if necessary, we may assume that
\[
  \zeta+\zeta^{-1}=\varphi-1.
\]
Put
\[
  \theta=\zeta-\zeta^{-1}.
\]
Then
\[
  \theta^2
  =
  (\zeta+\zeta^{-1})^2-4
  =
  (\varphi-1)^2-4
  =
  -(\varphi+2).
\]
Because $p\equiv1\pmod4$, $-1$ is a square in $\mathbb F_p$. Hence
$\varphi+2$ is a square in $\mathbb F_p$ if and only if $\theta^2$ is a square
in $\mathbb F_p$.

If $p\equiv1\pmod5$, then $\zeta\in\mathbb F_p$, hence
$\theta\in\mathbb F_p$, and therefore $\theta^2$, and so also $\varphi+2$,
is a square.

If $p\equiv4\pmod5$, then Frobenius acts on $\zeta$ by
\[
  \zeta^p=\zeta^{-1}.
\]
Consequently
\[
  \theta^p=(\zeta-\zeta^{-1})^p=\zeta^{-1}-\zeta=-\theta.
\]
Thus $\theta\notin\mathbb F_p$. Since the only square roots of $\theta^2$ are
$\pm\theta$, the element $\theta^2$ cannot be a square in $\mathbb F_p$.
Because $-1$ is a square, $\varphi+2=-\theta^2$ is also a nonsquare.

Therefore
\[
  \chi(\varphi+2)=
  \begin{cases}
    1,&p\equiv1\pmod5,\\
    -1,&p\equiv4\pmod5,
  \end{cases}
\]
which is precisely $\delta_p$.
\end{proof}

% =====================================================================
\section{The split case}\label{sec:split}
% =====================================================================

\begin{proposition}[Jacobsthal reduction, split]\label{prop:split-red}
Let $\varepsilon_p=1$ and $i_p=2$. Then
\[
  S(p)=\begin{cases}\qleg5p\cdot\tfrac12\Jac(-1), & p\equiv1\pmod8,\\[2pt] -\qleg5p, & p\equiv5\pmod8.\end{cases}
\]
\end{proposition}
\begin{proof}
By Lemma~\ref{lem:index-parity}, $p\equiv1\pmod4$, and $z=(p-1)/2$.

\emph{Step 1.} Put $g=\iota\varphi$ and $y_j=g^j$. Since
$\psi^j=(-1)^j\varphi^{-j}$,
$F_j=\varphi^{-j}\bigl(\varphi^{2j}-(-1)^j\bigr)/(\varphi-\psi)$. Now
$y_j^2=(-1)^j\varphi^{2j}$, so $\varphi^{2j}-(-1)^j=(-1)^j(y_j^2-1)$, and
$\chi(\varphi^{-j})=\chi(\varphi^{j})=\chi(\iota)^j\chi(y_j)$. Therefore
\[
  \chi(F_j)=\qleg5p\,\leg2p^{j} f(y_j),\qquad f(y):=\chi(y)\chi(y^2-1).
\]

\emph{Step 2.} Since $g^2=\eta$ has order $(p-1)/2$ and this order is even,
$g$ has order $p-1$. Indeed, if $g$ had order $(p-1)/2$, then $g^2$ would
have order $(p-1)/4$, a contradiction. Thus $g$ is a primitive root.
Consequently $\{y_j:1\le j\le z\}$ is a half-system of $\mathbb F_p^\times$
modulo $\pm1$, because $g^{j+z}=-g^j$. Also $f$ is even, since $\chi(-1)=1$.

\emph{Step 3.} If $p\equiv1\pmod 8$, then $\leg2p=1$, so
$S(p)=\qleg5p\sum_{j=1}^{z}f(y_j)=\qleg5p\cdot\tfrac12\sum_{y\neq0}f(y)
=\qleg5p\cdot\tfrac12\Jac(-1)$.

If $p\equiv5\pmod8$, then $\leg2p^{j}=(-1)^j=\chi(g^j)$, so
$S(p)=\qleg5p\sum_{j}\chi(y_j^2-1)=\qleg5p\cdot\tfrac12\sum_{y\ne0}\chi(y^2-1)$.
Since
\[
  \sum_{y\in\Fp}\chi(y^2-1)=-1
\]
and $\chi(-1)=1$, we have
\[
  \sum_{y\ne0}\chi(y^2-1)=-2.
\]
Consequently
\[
  S(p)=\qleg5p\cdot\tfrac12(-2)=-\qleg5p.
\]
\end{proof}

\begin{proof}[Proof of Theorem~\ref{thm:A}]
By Lemma~\ref{lem:quartic-unit},
\[
  \qleg5p\,\chi(\varphi)=\delta_p.
\]
Since $\chi(\varphi)^2=1$, this gives
\[
  \qleg5p=\delta_p\chi(\varphi).
\]
By Lemma~\ref{lem:split-index2-chi},
\[
  \chi(\varphi)=-\leg2p.
\]
Therefore
\[
  \qleg5p=-\delta_p\leg2p.
\]

If $p\equiv1\pmod8$, then $\leg2p=1$, so $\qleg5p=-\delta_p$.
By Proposition~\ref{prop:split-red} and Lemma~\ref{lem:Jminus1},
\[
  S(p)=\qleg5p\cdot\frac12\Jac(-1)
      =(-\delta_p)(-a^{*})
      =\delta_p a^{*}.
\]

If $p\equiv5\pmod8$, then $\leg2p=-1$, so $\qleg5p=\delta_p$.
By Proposition~\ref{prop:split-red},
\[
  S(p)=-\qleg5p=-\delta_p.
\]
\end{proof}

\begin{corollary}\label{cor:quartic-five}
If $\varepsilon_p=1$ and $i_p=2$, then $\qleg5p=-\delta_p\leg2p$.
Equivalently, $5$ is a quartic residue modulo such $p$ if and only if
$p\equiv9$ or $21\pmod{40}$.
\end{corollary}
\begin{proof}
The first identity was established in the proof of Theorem~\ref{thm:A}; the
residue classes follow by inspection of the four cases $p\bmod40$.
\end{proof}

\begin{remark}[Compatibility with a criterion of E.~Lehmer]
Emma Lehmer~\cite{ELehmer1958} proved that for primes $p\equiv1\pmod4$ with
$\leg5p=1$ one has $\qleg5p=1$ if and only if $5\mid b$. Combined with
Corollary~\ref{cor:quartic-five} this yields the consistency check that, for
split primes with $i_p=2$,
\[
  5\mid b\quad\Longleftrightarrow\quad p\equiv9\text{ or }21\pmod{40}.
\]
The equivalence has been verified for all such primes below $6\cdot10^4$
($550$ primes).
\end{remark}

\begin{remark}[Parity consistency]
For split $i_p=2$, $z(p)=(p-1)/2$ is even, so $S(p)$ is odd by
Lemma~\ref{lem:parity}, in accordance with the values $\delta_p a^{*}$ and $-\delta_p$.
For inert $i_p=2$, $z(p)$ is odd, so $S(p)$ is even, in accordance with
$\pm b$ and $0$.
\end{remark}

\begin{remark}
In terms of residue classes: $S(p)=a^{*}$ for $p\equiv1\pmod{40}$,
$S(p)=-a^{*}$ for $p\equiv9\pmod{40}$, $S(p)=-1$ for $p\equiv21\pmod{40}$, and
$S(p)=1$ for $p\equiv29\pmod{40}$. Compare the full-rank law $S(p)=\delta_p$ of
\cite{Ghandali2026Preprint}. For $i_p=2$ one has
\[
  S(p)\equiv \delta_p\leg2p \equiv -\qleg5p \pmod4.
\]
A uniform statement covering all indices remains open.
\end{remark}

% =====================================================================
\section{The inert case}\label{sec:inert}
% =====================================================================

Let $\varepsilon_p=-1$. Let $N:\Fpp^{\times}\to\Fp^{\times}$ be the norm and
$T=\ker N$, a cyclic group of order $p+1$. Put $\sqrt5:=\varphi-\psi$; then
$(\sqrt5)^p=-\sqrt5$.

\begin{lemma}[Cayley parametrisation]\label{lem:cayley}
The map $c(t)=(t+\sqrt5)/(t-\sqrt5)$, $c(\infty)=1$, is a bijection
$\mathbb P^1(\Fp)\to T$ with $c(0)=-1$ and $c(5/t)=-c(t)$. For $t\in\Fp^{\times}$
\[
  \frac{c(t)-c(t)^{-1}}{\sqrt5}=\frac{4t}{t^2-5}\in\Fp^{\times}.
\]
\end{lemma}
\begin{proof}
Standard: $c(t)^p=c(t)^{-1}$, and the inverse is $y\mapsto\sqrt5\,(y+1)/(y-1)$,
which is Frobenius-fixed. Since $5$ is a non-residue, $t^2\ne5$. The identities
follow by direct computation.
\end{proof}

\begin{proposition}[Jacobsthal reduction, inert]\label{prop:inert-red}
Let $\varepsilon_p=-1$ and $i_p=2$. If $p\equiv1\pmod 8$, then
\[
  S(p)=\frac12\Jac(-5)=\frac12\Jac(5).
\]
\end{proposition}

\begin{proof}
Here $p\equiv1\pmod4$ by Lemma~\ref{lem:index-parity}, and
\[
  z=z(p)=\frac{p+1}{2}
\]
is odd. Put
\[
  g=\iota\varphi.
\]
Then
\[
  N(g)=N(\iota)N(\varphi)=(-1)(-1)=1,
\]
so $g\in T$, and
\[
  g^2=\eta
\]
has order $z$.

With $y_j=g^j$ one has the identities
\[
  (-\iota)^j y_j=\varphi^j,\qquad
  (-\iota)^j y_j^{-1}=(-1)^j\varphi^{-j}=\psi^j,
\]
and therefore
\[
  F_j=(-\iota)^{j}\,\frac{y_j-y_j^{-1}}{\sqrt5}.
\]
(For $y\in T$, the quotient $(y-y^{-1})/\sqrt5$ lies in $\Fp$, because
Frobenius sends $y-y^{-1}$ to $-(y-y^{-1})$ and sends $\sqrt5$ to $-\sqrt5$.)
Consequently
\[
  \chi(F_j)=\leg2p^{j}\,
  \chi\!\left(\frac{y_j-y_j^{-1}}{\sqrt5}\right).
\]
Since $p\equiv1\pmod8$, we have $\leg2p=1$, so this becomes
\[
  \chi(F_j)=
  \chi\!\left(\frac{y_j-y_j^{-1}}{\sqrt5}\right).
\]

We next record explicitly how the powers of $g$ sample the norm-one torus.
Since $g^2=\eta$ has order $z=(p+1)/2$, the order of $g$ is either $z$ or
$2z=p+1$. In both cases the set
\[
  Y=\{g^j:1\le j\le z\}
\]
contains exactly one element from each pair $\{y,-y\}$ in $T$.

Indeed, if $\operatorname{ord}(g)=2z$, then $g^z=-1$, and therefore
\[
  g^{j+z}=-g^j.
\]
Thus the first $z$ powers give one representative from each pair
$\{y,-y\}$.

If instead $\operatorname{ord}(g)=z$, then $\langle g\rangle$ is the unique
subgroup of $T$ of order $z$. Since $z$ is odd, this subgroup does not contain
$-1$. Hence
\[
  T=\langle g\rangle\sqcup(-\langle g\rangle),
\]
and again $Y=\langle g\rangle$ contains exactly one element from each pair
$\{y,-y\}$.

Now use the Cayley parametrisation of Lemma~\ref{lem:cayley}. Write
\[
  y=c(t)=\frac{t+\sqrt5}{t-\sqrt5}.
\]
For $t\in\mathbb F_p^{\times}$, Lemma~\ref{lem:cayley} gives
\[
  \frac{c(t)-c(t)^{-1}}{\sqrt5}
  =
  \frac{4t}{t^2-5}.
\]
Therefore
\[
  \chi\!\left(\frac{c(t)-c(t)^{-1}}{\sqrt5}\right)
  =
  \chi(t)\chi(t^2-5).
\]
Put
\[
  f(t):=\chi(t)\chi(t^2-5).
\]

The involution $y\mapsto -y$ on $T$ corresponds, under the Cayley
parametrisation, to
\[
  t\mapsto \frac5t.
\]
(Alternatively, without Cayley: $(-y)-(-y)^{-1}=-(y-y^{-1})$, so the
character is invariant under $y\mapsto-y$ as soon as $\chi(-1)=1$, i.e.,
already under the weaker hypothesis $p\equiv1\pmod4$.)
Moreover,
\[
  f\!\left(\frac5t\right)=f(t).
\]
Indeed, a direct computation yields
\[
\begin{aligned}
 f\!\left(\frac5t\right)
 &= \chi\!\left(\frac5t\right)
    \chi\!\left(\frac{25}{t^2}-5\right)  \\
 &= \chi(5)\chi(t)\,
    \chi\!\left(\frac{5(5-t^2)}{t^2}\right) \\
 &= \chi(5)^2\chi(t)\chi(5-t^2) \\
 &= \chi(t)\chi(-(t^2-5)) \\
 &= \chi(t)\chi(t^2-5) \\
 &= f(t),
\end{aligned}
\]
where the last step uses only $\chi(-1)=1$ (i.e., $p\equiv1\pmod4$).

The pair $\{1,-1\}\subset T$ corresponds to $t=\infty$ and $t=0$.
Both give zero contribution: for $y=1$ one has $y-y^{-1}=0$, and for
$y=-1$ one again has $y-y^{-1}=0$. Hence summing over one representative
from each pair $\{y,-y\}$ is the same as taking one half of the sum over
all nonzero finite parameters $t\in\mathbb F_p^\times$. Therefore
\[
  S(p)
  =
  \frac12\sum_{t\in\mathbb F_p^\times}
  \chi(t)\chi(t^2-5).
\]
The term $t=0$ is zero, so
\[
  S(p)
  =
  \frac12\sum_{t\in\mathbb F_p}
  \chi(t)\chi(t^2-5)
  =
  \frac12\Jac(-5).
\]

Finally, substituting $x=\iota u$ gives
\[
\begin{aligned}
 \Jac(-5)
 &= \sum_{x\in\mathbb F_p}\chi(x)\chi(x^2-5)\\
 &= \sum_{u\in\mathbb F_p}\chi(\iota u)\chi(-u^2-5)\\
 &= \chi(\iota)\chi(-1)
    \sum_{u\in\mathbb F_p}\chi(u)\chi(u^2+5)\\
 &= \leg2p\,\Jac(5).
\end{aligned}
\]
Since $p\equiv1\pmod8$, $\leg2p=1$, and hence
\[
  \Jac(-5)=\Jac(5).
\]
This proves the proposition.
\end{proof}

\begin{definition}[{Quartic residue symbol in $\mathbb Z[i]$}]\label{def:quartic-symbol}
Let $\pi$ be a primary Gaussian prime of odd norm
\[
  N(\pi)=p,
\]
and let $\alpha\in\mathbb Z[i]$ be coprime to $\pi$.
The quartic residue symbol
\[
  \biggl(\frac{\alpha}{\pi}\biggr)_{\!4}\in\{1,-1,i,-i\}
\]
is the unique fourth root of unity satisfying the Euler criterion
\[
  \alpha^{(p-1)/4}
  \equiv
  \biggl(\frac{\alpha}{\pi}\biggr)_{\!4}
  \pmod{\pi}.
\]
An odd Gaussian integer is called primary if it is congruent to $1$
modulo $2+2i$. We use the standard quartic reciprocity law for
primary Gaussian integers; see
Ireland--Rosen~\cite[Chapter~9, §§8--9, pp.~121--126]{IR}.
\end{definition}

\begin{lemma}[Sign via biquadratic reciprocity]\label{lem:sign5}
Let $p\equiv1\pmod8$ with $\leg5p=-1$, and write
\[
  p=(a^{*})^2+b^2,
  \qquad
  a^{*}\equiv1\pmod4,\qquad b>0.
\]
Put
\[
  \pi=a^{*}+bi.
\]
With the convention of Definition~\ref{def:quartic-symbol}, one has
$5\nmid a^{*}b$,
\[
  b\equiv\pm a^{*}\pmod5,
\]
and
\[
  \biggl(\frac{5}{\pi}\biggr)_{\!4}
  =
  \begin{cases}
    i,  & b\equiv a^{*}\pmod5,\\
    -i, & b\equiv -a^{*}\pmod5.
  \end{cases}
\]
\end{lemma}

\begin{proof}
We use Definition~\ref{def:quartic-symbol} throughout. Since
$\leg5p=-1$, one has
\[
  p\equiv2\ \text{or}\ 3\pmod5.
\]
The quadratic residues modulo~$5$ are $0,1,4$. From
\[
  p=(a^{*})^2+b^2
\]
it follows that neither $a^{*}$ nor $b$ is divisible by~$5$, and the two
squares must be equal modulo~$5$: they are both~$1$ if $p\equiv2\pmod5$,
and both~$4$ if $p\equiv3\pmod5$. Hence
\[
  b^2\equiv(a^{*})^2\pmod5,
\]
and therefore
\[
  r:=\frac{b}{a^{*}}\equiv\pm1\pmod5.
\]

Write
\[
  5=\lambda\bar\lambda,\qquad
  \lambda=-1+2i,\qquad
  \bar\lambda=-1-2i.
\]
Both $\lambda$ and $\bar\lambda$ are primary. Also $\pi$ is primary:
since $p\equiv1\pmod8$, one has $b\equiv0\pmod4$ and
$a^{*}\equiv1\pmod4$, so $\pi\equiv1\pmod{2+2i}$.

For relatively prime primary Gaussian primes $\alpha$ and $\beta$,
quartic reciprocity gives
\[
  \biggl(\frac{\alpha}{\beta}\biggr)_{\!4}
  =
  (-1)^{\frac{N(\alpha)-1}{4}\frac{N(\beta)-1}{4}}
  \biggl(\frac{\beta}{\alpha}\biggr)_{\!4};
\]
see Ireland--Rosen~\cite[Chapter~9, §9, pp.~123--126]{IR}.
Since
\[
  N(\lambda)=N(\bar\lambda)=5
\]
and
\[
  \frac{p-1}{4}
\]
is even, the sign factor is trivial. Therefore
\[
  \biggl(\frac{\lambda}{\pi}\biggr)_{\!4}
  =
  \biggl(\frac{\pi}{\lambda}\biggr)_{\!4},
  \qquad
  \biggl(\frac{\bar\lambda}{\pi}\biggr)_{\!4}
  =
  \biggl(\frac{\pi}{\bar\lambda}\biggr)_{\!4}.
\]
Hence
\[
  \biggl(\frac{5}{\pi}\biggr)_{\!4}
  =
  \biggl(\frac{\pi}{\lambda}\biggr)_{\!4}
  \biggl(\frac{\pi}{\bar\lambda}\biggr)_{\!4}.
\]

Under the isomorphism $\mathbb Z[i]/(\lambda)\simeq\mathbb F_5$ one has
$i\mapsto3$, while under
$\mathbb Z[i]/(\bar\lambda)\simeq\mathbb F_5$ one has $i\mapsto2=3^{-1}$.
Let $\ell$ denote the discrete logarithm to base~$3$ on $\mathbb F_5^\times$.
Modulo $\lambda$,
\[
  \pi=a^{*}+bi\equiv a^{*}(1+3r),
\]
and modulo $\bar\lambda$,
\[
  \pi=a^{*}+bi\equiv a^{*}(1+2r).
\]
The factor $a^{*}$ cancels in the product.  Indeed, writing
$a^{*}\equiv 3^{k}\pmod5$, its contribution modulo~$\lambda$ is
$i^{k}$ (under $i\mapsto3$), while its contribution
modulo~$\bar\lambda$ is $i^{-k}$ (under $i\mapsto3^{-1}$); the product
of these contributions is~$1$.  Therefore
\[
  \biggl(\frac{5}{\pi}\biggr)_{\!4}
  =
  i^{\,\ell((1+3r)/(1+2r))}.
\]
If $r=1$, then
\[
  \frac{1+3r}{1+2r}=\frac4{3}\equiv3\pmod5,
\]
so the exponent is~$1$ and the symbol is~$i$. If $r=-1$, then
\[
  \frac{1+3r}{1+2r}=\frac{-2}{-1}\equiv2\pmod5,
\]
so the exponent is~$3$ and the symbol is $i^3=-i$. This proves the claim.
\end{proof}

\begin{proof}[Proof of Theorem~\ref{thm:B}]
By Lemma~\ref{lem:index-parity}, $p\equiv1\pmod4$ and $z(p)=(p+1)/2$.
If $p\equiv5\pmod8$, then $S(p)=0$ by Corollary~\ref{cor:vanishing}(a).

Now assume $p\equiv1\pmod8$. Proposition~\ref{prop:inert-red} gives
$S(p)=\tfrac12\Jac(5)$. Since $\varepsilon_p=-1$, $5$ is a quadratic
non-residue modulo $p$. Choose any nonzero quadratic residue $r\bmod p$.
By Lemma~\ref{lem:jacobsthal}, $\tfrac12\Jac(r)$ is odd,
$\tfrac12\Jac(5)$ is even, and
\[
  \bigl(\tfrac12\Jac(r)\bigr)^2+\bigl(\tfrac12\Jac(5)\bigr)^2=p.
\]
The odd component of a representation of $p$ as a sum of two squares is
$|a^{*}|$, and the even component is $b$. Hence
$|\tfrac12\Jac(5)|=b$.

It remains to determine the sign. By Lemma~\ref{lem:jac-cong},
\[
  \tfrac12\Jac(5)\equiv -a^{*}\,5^{(p-1)/4}\pmod p.
\]
Let
\[
  \pi=a^{*}+bi.
\]
Under the isomorphism
\[
  \mathbb Z[i]/(\pi)\simeq\mathbb F_p
\]
one has
\[
  i\equiv -\frac{a^{*}}{b}\pmod\pi.
\]
Moreover,
\[
  5^{(p-1)/4}\equiv
  \biggl(\frac{5}{\pi}\biggr)_{\!4}
  \quad\text{in }\mathbb Z[i]/(\pi).
\]

If $b\equiv a^{*}\pmod5$, Lemma~\ref{lem:sign5} gives
\[
  \biggl(\frac{5}{\pi}\biggr)_{\!4}=i.
\]
Hence, in $\mathbb F_p$,
\[
  5^{(p-1)/4}\equiv -\frac{a^{*}}{b}.
\]
Therefore
\[
  \tfrac12\Jac(5)
  \equiv
  -a^{*}\Bigl(-\frac{a^{*}}{b}\Bigr)
  =
  \frac{(a^{*})^2}{b}
  \equiv
  -b
  \pmod p,
\]
because $(a^{*})^2+b^2=p\equiv0\pmod p$.

If $b\equiv -a^{*}\pmod5$, Lemma~\ref{lem:sign5} gives
\[
  \biggl(\frac{5}{\pi}\biggr)_{\!4}=-i.
\]
Hence
\[
  5^{(p-1)/4}\equiv \frac{a^{*}}{b}\pmod p,
\]
and therefore
\[
  \tfrac12\Jac(5)
  \equiv
  -a^{*}\frac{a^{*}}{b}
  =
  -\frac{(a^{*})^2}{b}
  \equiv
  b
  \pmod p.
\]

Since $\tfrac12\Jac(5)=\pm b$ and $b\not\equiv -b\pmod p$, these congruences
determine the sign:
\[
  \tfrac12\Jac(5)
  =
  \begin{cases}
    -b,& b\equiv a^{*}\pmod5,\\
    \phantom{-}b,& b\equiv -a^{*}\pmod5.
  \end{cases}
\]
Equivalently, $S(p)$ is the unique element of $\{\pm b\}$ congruent to
$-a^{*}$ modulo~$5$.
\end{proof}

\begin{remark}
Split index two thus realises the \emph{residue} Jacobsthal component ($\pm a^{*}$),
and inert index two realises the \emph{non-residue} component ($\pm b$).
\end{remark}

\begin{remark}[Why parameter $5$ rather than $10$]
For $p\equiv1\pmod8$ write $2=s^2$. Then
$\Jac(10)=\Jac(5s^2)=\chi(s)\Jac(5)=\qleg2p\Jac(5)$.
Hence $|\Jac(\pm10)|=2|S(p)|$, but the sign is twisted by the quartic symbol
$\qleg2p$.
\end{remark}
% =====================================================================
\section{Numerical evidence}\label{sec:numerics}
% =====================================================================

The companion verifier \texttt{ghandali\_verify\_index2.py}
(version~1.0.1), archived with its machine-readable outputs and run log
in~\cite{ghandaliIndex2Archive}, was run with
\[
  N=2\cdot 10^{6},
  \qquad
  \texttt{sum\_limit}=2\cdot 10^{6},
  \qquad
  \texttt{jac\_limit}=2\cdot 10^{6},
\]
computing direct Fibonacci character sums for rank indices \(2\) and
\(4\).  The run formed a rank inventory of \(148{,}931\) primes
\(p\le N\) with \(p\nmid10\); all theorem-level cases reported in the
table below satisfy the standing hypothesis \(p>5\).  It selected
\(52{,}477\) direct \(S(p)\)-records and returned zero mismatches in
the theorem-level identities checked by the verifier.

\begin{center}
\small
\begin{tabular}{@{}llp{1.6cm}p{1.8cm}l@{}}
\toprule
Case & Prediction & Primes & Agreement & Proved in\\
\midrule
inert, \(i_p=2\), \(p\equiv1\pmod8\)
  & \(S\in\{\pm b\}\), \(S\equiv-a^{*}\pmod5\)
  & 14591 & 14591/14591 & Thm.~\ref{thm:B}\\
inert, \(i_p=2\), \(p\equiv5\pmod8\)
  & \(S=0\)
  & 14696 & 14696/14696 & Cor.~\ref{cor:vanishing}\\
split, \(i_p=2\), \(p\equiv1\pmod8\)
  & \(S=\delta_p a^{*}\)
  & 6551 & 6551/6551 & Thm.~\ref{thm:A}\\
split, \(i_p=2\), \(p\equiv5\pmod8\)
  & \(S=-\delta_p\)
  & 6583 & 6583/6583 & Thm.~\ref{thm:A}\\
split, \(i_p=4\), \(p\equiv5\pmod8\)
  & \(S=0\)
  & 6697 & 6697/6697 & Cor.~\ref{cor:vanishing}\\
split, \(i_p=4\), \(p\equiv1\pmod8\)
  & open; \(S\) odd
  & 3359 & --- & ---\\
\bottomrule
\end{tabular}
\end{center}

The same run also checked the rank identity \(z(p)i_p=n_p\) and the
index-parity criterion on all \(148{,}931\) primes in the rank
inventory.  The identity
\[
  \chi(F_{z(p)+1})=\left(\frac2p\right)
\]
for odd \(z(p)\) was checked on \(49{,}618\) primes, and the reflection
criteria used in Corollary~\ref{cor:vanishing} were checked with zero
mismatches.  The split and inert Jacobsthal reductions were checked
directly throughout the same range wherever the corresponding direct
Jacobsthal sums were computed; in particular,
\[
  2S(p)=\Jac(5)
\]
was checked for all \(14{,}591\) inert index-two primes with
\(p\equiv1\pmod8\) below \(2\cdot10^{6}\).  The same run verified the
signed law \(S(p)\equiv-a^{*}\pmod5\) on those primes with zero
mismatches.

Claims of zero mismatches refer to the archived run deposited
in~\cite{ghandaliIndex2Archive}, which contains the verifier source code
(version~1.0.1), machine-readable CSV and JSON output files, the run log
with the parameters recorded above, and checksums.  The present paper
records only the mathematical summary of these computations.

% =====================================================================
\section{Conditional density context}\label{sec:context}
% =====================================================================

This section is included only as conditional context under the
generalized Riemann hypothesis and is not used in the proofs of
Theorems~\ref{thm:A} and~\ref{thm:B}. Those theorems are unconditional.

\subsection{A conditional density corollary for inert index two}\label{sec:density}

The computations in Section~\ref{sec:numerics} suggest the two index-two
density constants
\[
  \frac{9A}{38}\qquad\text{in the split branch}
\]
and
\[
  \frac{10A}{19}\qquad\text{in the inert branch},
\]
where \(A\) denotes Artin's constant.  The split value is recorded here
as a conjectural density suggested by the data. In contrast, the inert value follows, as a Dirichlet density, conditionally from the
GRH-dependent density theorems of Roskam~\cite{Roskam2000} and
Chen~\cite{Chen2002}, in the general Artin-type framework of
Lenstra~\cite{Lenstra1977}, once their maximal-order statements are
translated into the present rank-index notation; the derivation is
self-contained in Corollary~\ref{cor:inert-index-two-density} and does
not rely on~\cite{Ghandali2026Preprint}.

\paragraph{The rank index as a group index.}
Let \(p\nmid 10\), and write
\[
  \eta=\frac{\varphi}{\psi}=-\varphi^{2}.
\]
If \(\leg{5}{p}=1\), put \(G_p=\mathbb F_p^{\times}\).  If
\(\leg{5}{p}=-1\), let
\[
  G_p=T(\mathbb F_p)
  =
  \ker\!\left(
    N:\mathbb F_{p^{2}}^{\times}\to\mathbb F_p^{\times}
  \right)
\]
be the norm-one torus.  In both cases \(G_p\) is cyclic of order
\[
  n_p=p-\leg{5}{p},
\]
and \(\eta\in G_p\).  Since
\[
  F_n=\frac{\varphi^{n}-\psi^{n}}{\sqrt5},
\]
one has \(F_n\equiv0\pmod p\) if and only if \(\eta^n=1\).  Hence
\[
  z(p)=\operatorname{ord}_{G_p}(\eta),
  \qquad
  i_p=[G_p:\langle\eta\rangle].
\]
Thus full rank \(i_p=1\) means that \(\eta\) is a primitive point of
the one-dimensional torus attached to \(\mathbb Q(\sqrt5)\), while
rank index \(t\) means that \(\eta\) generates the unique subgroup of
index \(t\).  This is the torus-theoretic form of the generalized
Artin problem in the sense of Lenstra~\cite{Lenstra1977}, and is the
setting of Chen's theorem on primitive points of one-dimensional
tori~\cite{Chen2002}.
\begin{lemma}[2-adic dictionary, inert branch]\label{lem:2adic}
Let \(p\neq2,5\) be inert in \(\mathbb Q(\sqrt5)\), and put
\[
  d=\operatorname{ord}(\bar\varphi^{\,2}),
\]
where \(\bar\varphi\) denotes the image of \(\varphi\) in
\(\mathbb F_{p^2}^{\times}\).  Then:
\begin{enumerate}[label=\textup{(\roman*)}]
\item
\[
  \bar\varphi^{\,p+1}=N(\bar\varphi)=-1.
\]
Consequently \(d\mid p+1\), the quotient \((p+1)/d\) is odd, and
\[
  \bar\eta=(\bar\varphi^{\,2})^{(p+3)/2}.
\]

\item
If \(p\equiv3\pmod4\), then
\[
  \gcd\!\left(d,\frac{p+3}{2}\right)=1,
\]
so
\[
  z(p)=d,\qquad i_p=\frac{p+1}{d},
\]
and \(i_p\) is odd.

\item
If \(p\equiv1\pmod4\), then \(d=2u\) with \(u\) odd and
\[
  \gcd\!\left(d,\frac{p+3}{2}\right)=2.
\]
Therefore
\[
  z(p)=u=\frac d2,\qquad
  i_p=\frac{p+1}{u}=\frac{2(p+1)}{d}\equiv2\pmod4.
\]
\end{enumerate}
In particular,
\begin{equation}\label{eq:dictionary}
\{\,p\text{ inert}:\ i_p\in\{1,2\}\,\}
=
\{\,p\text{ inert}:\ d=p+1\,\}.
\end{equation}
The set on the right is the inert maximal-order set for the fundamental
unit \(\varphi\), equivalently the set denoted \(S^{-}\) in Roskam's
specialization to \(\mathbb Q(\sqrt5)\)~\cite{Roskam2000}.
\end{lemma}

\begin{proof}
Since \(p\) is inert in \(\mathbb Q(\sqrt5)\), Frobenius interchanges
\(\bar\varphi\) and \(\bar\psi\).  Hence
\[
  \bar\varphi^{\,p+1}
  =
  \bar\varphi\,\bar\varphi^{\,p}
  =
  \bar\varphi\,\bar\psi
  =
  -1.
\]
Let \(m=\operatorname{ord}(\bar\varphi)\).  Then
\(m\mid2(p+1)\), while \(m\nmid p+1\), because
\(\bar\varphi^{\,p+1}=-1\).  Therefore
\[
  d=\operatorname{ord}(\bar\varphi^{\,2})
\]
divides \(p+1\), but \(d\nmid(p+1)/2\).  Equivalently,
\((p+1)/d\) is odd.  Moreover
\[
  -1
  =
  \bar\varphi^{\,p+1}
  =
  (\bar\varphi^{\,2})^{(p+1)/2},
\]
and hence
\[
  \bar\eta=-\bar\varphi^{\,2}
  =
  (\bar\varphi^{\,2})^{(p+3)/2}.
\]
This proves (i).

For (ii), write \(p=4k+3\).  Then \((p+3)/2=2k+3\) is odd.  Any common
divisor of \(d\) and \((p+3)/2\) divides
\[
  2\cdot\frac{p+3}{2}-(p+1)=2.
\]
Since \((p+3)/2\) is odd, this common divisor must be \(1\).  Therefore
\[
  z(p)
  =
  \operatorname{ord}(\bar\eta)
  =
  \frac{d}{\gcd(d,(p+3)/2)}
  =
  d,
\]
and so
\[
  i_p=\frac{p+1}{z(p)}=\frac{p+1}{d},
\]
which is odd by (i).

For (iii), write \(p=4k+1\).  Then \(v_2(p+1)=1\).  Since
\((p+1)/d\) is odd, we have \(v_2(d)=1\), so \(d=2u\) with \(u\) odd.
Also
\[
  \frac{p+3}{2}=\frac{p+1}{2}+1
\]
is even, and
\[
  \gcd\!\left(u,\frac{p+3}{2}\right)
  \mid
  \gcd\!\left(\frac{p+1}{2},\frac{p+3}{2}\right)
  =
  1.
\]
Thus
\[
  \gcd\!\left(d,\frac{p+3}{2}\right)=2.
\]
It follows that
\[
  z(p)=\frac d2=u,
  \qquad
  i_p=\frac{p+1}{u}=\frac{2(p+1)}{d}.
\]
Since \((p+1)/d\) is odd, this gives \(i_p\equiv2\pmod4\).

Finally, if \(p\equiv3\pmod4\), then \(i_p\) is odd and
\(i_p=1\) is equivalent to \(d=p+1\).  If \(p\equiv1\pmod4\), then
\(i_p\equiv2\pmod4\), and \(i_p=2\) is again equivalent to
\(d=p+1\).  This proves \eqref{eq:dictionary}.  The final identification with Roskam's inert maximal-order set is
exactly the condition that \(\bar\varphi^{\,2}\) generate the norm-one
group of order \(p+1\), as in Roskam's quadratic analogue of Artin's
conjecture~\cite{Roskam2000}.
\end{proof}
\begin{corollary}[GRH; Roskam--Chen]
\label{cor:inert-index-two-density}
Assume the generalized Riemann hypothesis for the Kummer extensions
occurring in the density theorems of Roskam~\cite{Roskam2000} and
Chen~\cite{Chen2002}. Then the set of primes $p>5$ such that
\[
\leg{5}{p}=-1,
\qquad
i_p=2
\]
has Dirichlet density
\[
\frac{10A}{19}
\]
among all rational primes, where $A$ denotes Artin's constant. This
is a specialization of the cited conditional density theorems into
the present rank-index notation; it is not an independent density
theorem of this paper and is not used in the proofs of
Theorems~\ref{thm:A} and~\ref{thm:B}.
\end{corollary}

\begin{proof}
By Lemma~\ref{lem:2adic}, the inert primes with $i_p\in\{1,2\}$ are
precisely the inert primes for which
\[
\operatorname{ord}(\bar\varphi^{\,2})=p+1.
\]
This is Roskam's inert maximal-order set $S^-$ for the fundamental
unit $\varphi$ of $\mathbb Q(\sqrt5)$. Since $N(\varphi)=-1$ and
$D=5$ is odd, Roskam's GRH-conditional Theorem~2(b)
\cite{Roskam2000} gives the \emph{natural} density
\[
\delta_{\mathrm{nat}}(S^-)=A\bigl(1-E(5)\bigr)
=A\Bigl(1+\frac1{19}\Bigr)=\frac{20A}{19},
\]
where, for square-free $a$,
\[
E(a)=\mu(a)\prod_{q\mid a}\frac{1}{q(q-1)-1},
\qquad
E(5)=-\frac{1}{5\cdot4-1}=-\frac1{19}.
\]
The erratum to \cite{Roskam2000} corrects Lemmas~6--7 and the
formulas of Theorem~2(a) for the case $N(\varepsilon)=+1$ and for
$\delta(S^+)$ when $2\mid D$; it does not alter the formula
$\delta(S^-)=A(1-E(D))$ used here for $N(\varepsilon)=-1$. A set
with natural density has the same Dirichlet density, so
\[
\delta_D(S^-)=\frac{20A}{19}.
\]

We now compute the Dirichlet density of the inert full-rank set
directly from Chen's Theorem~2.1 \cite{Chen2002}, applied to
\[
\eta=-\varphi^2\in K=\mathbb Q(\sqrt5).
\]
The element $\eta$ is nontorsion with $N_{K/\mathbb Q}(\eta)=1$ and
is not a proper power in $K^\times$: it is negative under both real
embeddings of $K$, so it has no square root in the real field $K$;
and an equality $\eta=(\pm\varphi^m)^q$ with $q$ an odd prime would
force $mq=2$, which is impossible. Hence the exponent $s$ in Chen's
notation equals $1$.

For square-free $n$, Chen's field $k_n$ equals $K$ if $n$ is odd and
$K(\sqrt\eta)=K(i)$ if $n$ is even. Consequently the intersection
degree
\[
[\,k_n\cap\mathbb Q(\mu_n):\mathbb Q\,]
=
\begin{cases}
2,&5\mid n,\\
1,&5\nmid n,
\end{cases}
\]
because $K=\mathbb Q(\sqrt5)$ is contained in $\mathbb Q(\mu_n)$
exactly when $5\mid n$, while a square-free $n$ is never divisible
by $4$, so $\mathbb Q(\mu_n)$ cannot contain $i$. Chen's
Lemma~1.6 \cite{Chen2002} therefore gives the Kummer degrees
\[
d_n=[E_n:\mathbb Q]=
\begin{cases}
n\varphi(n),&5\mid n,\\[2pt]
2n\varphi(n),&5\nmid n.
\end{cases}
\]
For the inert obstruction coefficients, Chen's Lemma~1.7(a)
\cite{Chen2002} gives $c_n^-=0$ when $5\mid n$, since $K$ is real
and contained in $\mathbb Q(\mu_n)$; and $c_n^-=1$ when
$5\nmid n$, since then none of the three obstruction conditions of
that lemma applies: $K$ is not contained in $\mathbb Q(\mu_n)$, the
field $K'=\mathbb Q(i)$ is not contained in $\mathbb Q(\mu_n)$ for
square-free $n$, and $K''=\mathbb Q(\sqrt{-5})$ is imaginary.
Chen's formula
$\operatorname{den}(M_\eta^-)=\sum_n c_n^-\,\mu(n)/d_n$
\cite[Theorem~2.1]{Chen2002} therefore yields
\[
\delta_D(M_\eta^-)
=\sum_{5\nmid n}\frac{\mu(n)}{2n\varphi(n)}
=\frac12\prod_{q\ne5}\Bigl(1-\frac1{q(q-1)}\Bigr)
=\frac12\cdot\frac{A}{1-1/20}
=\frac{10A}{19}.
\]
Since $z(p)=\operatorname{ord}(\bar\eta)$ and, for inert $p$,
$i_p=(p+1)/z(p)$, the set $M_\eta^-$ is exactly the inert full-rank
set $\{p\text{ inert}: i_p=1\}$.

The sets $\{i_p=1\}$ and $\{i_p=2\}$ are disjoint
(Lemma~\ref{lem:2adic}(ii)--(iii): $i_p=1$ forces
$p\equiv3\pmod4$, while $i_p=2$ forces $p\equiv1\pmod4$), and their
union is $S^-$ up to finitely many excluded primes. Dirichlet
density is additive for finite disjoint decompositions, so
\[
\begin{aligned}
\delta_D\{\,p>5:\leg{5}{p}=-1,\ i_p=2\,\}
&=\delta_D(S^-)-\delta_D(M_\eta^-)\\
&=\frac{20A}{19}-\frac{10A}{19}
=\frac{10A}{19}.
\end{aligned}
\]
\end{proof}

\begin{remark}[Natural versus Dirichlet density]
\label{rem:natural-vs-dirichlet}
The corollary is stated for Dirichlet density, which is what the
above argument yields. Roskam's Theorem~2(b) in fact provides the
natural density of $S^-$, but Chen's Theorem~2.1 is stated for
Dirichlet density; subtracting a Dirichlet density from a natural
density does not, by itself, establish natural density for the
difference set. Promoting the statement to natural density would
require an effective Hooley-type argument (GRH-conditional Chebotarev
estimates with tail control) applied to the inert full-rank set
$\{i_p=1\}$ directly, along the lines of \cite{Murty1983}; the
convergence of the Euler product in the computation above is not by
itself such an estimate. The numerical data of
Section~\ref{sec:numerics} is of course consistent with either notion
of density. We record the Dirichlet-density statement, which is the
strongest that the cited references directly support.
\end{remark}

\begin{remark}[Numerical confirmation and the split branch]
\label{rem:split-status}
The Dirichlet-density value in Corollary~\ref{cor:inert-index-two-density} is also reflected numerically.
In the run up to \(2\cdot10^{6}\), there are
\[
  14{,}591+14{,}696=29{,}287
\]
inert primes with \(i_p=2\).  Writing \(C(N)=29{,}287\) for this count and
letting \(\pi^{*}(N)\) denote the number of primes \(p\le N\) with
\(p\nmid10\), one has
\[
  \frac{C(N)}{\pi^{*}(N)A}\approx0.525859,
\]
compared with the GRH-conditional prediction
\[
  \frac{10}{19}\approx0.526316.
\]
The relative deviation is about \(0.09\%\).

The split index-two density suggested by the data is
\[
  \frac{9A}{38}.
\]
In the same run, if \(C_{\mathrm{split}}(N)\) denotes the number of split
index-two primes up to \(N\), one has
\[
  \frac{C_{\mathrm{split}}(N)}{\pi^{*}(N)A}\approx0.235826,
\]
compared with
\[
  \frac{9}{38}\approx0.236842.
\]
Unlike the inert branch, the split index-two set does not appear to
admit the same direct dictionary with Roskam's split maximal-order set.
The obstruction is the \(2\)-primary Kummer behaviour described in
Remark~\ref{rem:kummer-halving}; see also the discussion of the
full-rank split branch in~\cite{Ghandali2026Preprint}.  Thus the split value is recorded
here as a conjectural density supported by the computation, not as a
theorem.  Its eventual proof should fall within the near-primitive-root
framework treated in
\cite{LenstraMoreeStevenhagen2014,Palenstijn2014,Kataoka2003}.
\end{remark}

\begin{remark}[Kummer degree when \(4\mid n\)]\label{rem:kummer-halving}
Write
\[
  K_n=\mathbb Q(\sqrt5,\zeta_n).
\]
Density computations of Artin type for the subgroup generated by
\(\eta\) depend on the Kummer degrees
\[
  [K_n(\eta^{1/n}):K_n].
\]
If \(4\mid n\), then \(i=\zeta_4\in K_n\), and
\[
  \eta=-\varphi^{2}=(i\varphi)^2
\]
is already a square in \(K_n\).  Consequently
\[
  [K_n(\eta^{1/n}):K_n]\le \frac n2.
\]
Thus any density computation which treats this degree as \(n\) without
correction is valid only away from the \(4\mid n\) levels.  This
\(2\)-primary obstruction is precisely why the split index-two density
requires a separate near-primitive-root analysis, rather than a naive
subtraction from Roskam's split maximal-order set.  Compare the
general Artin-type frameworks of Lenstra~\cite{Lenstra1977} and
Chen~\cite{Chen2002}.
\end{remark}

\paragraph{Numerical comparison.}
Table~\ref{tab:density} records the aggregate density ratios obtained
from the \(N=2\cdot10^{6}\) verification run.  The density rows are
included as context for the rank-index distribution; the finite
computation is not used as a proof of the GRH-conditional density
Corollary~\ref{cor:inert-index-two-density}.

\begin{table}[ht]
\centering
\small
\begin{tabular}{@{}llcc@{}}
\toprule
Quantity & Unit & Prediction & \(N=2\cdot10^{6}\) \\
\midrule
full rank, all & \(A\) & \(1\) & \(1.00128\) \\
full rank, split & \(A\) & \(9/19=0.47368\) & \(0.47415\) \\
full rank, inert & \(A\) & \(10/19=0.52632\) & \(0.52713\) \\
\midrule
\(i_p=2\), split & \(A\) & \(9/38=0.23684\) conjectural & \(0.23583\) \\
\(i_p=2\), inert & \(A\) & \(10/19=0.52632\) Dirichlet, conditional & \(0.52586\) \\
\bottomrule
\end{tabular}
\caption{Aggregate rank-index density ratios from the companion
verification run with \(N=2\cdot10^{6}\).  The density framework is the
GRH-conditional Artin--Lenstra--Chen--Roskam framework
\cite{Lenstra1977,Chen2002,Roskam2000}.  The inert index-two value
\(10A/19\) is the \emph{Dirichlet} density of
Corollary~\ref{cor:inert-index-two-density}
(see Remark~\ref{rem:natural-vs-dirichlet}).  The split index-two value
\(9A/38\) is recorded only as a conjectural density suggested by the
data.  The finite computation is not used as a proof of any density
statement.}
\label{tab:density}
\end{table}

\subsection{Framework and the series}

\begin{remark}[Framework for index densities]
The determination of densities at fixed rank index is an instance of
the generalized Artin problem in the sense of
Lenstra~\cite{Lenstra1977}.  Murty~\cite{Murty1983} supplies a
conditional analytic framework for Artin-type problems.  The
quadratic-field case for units is treated by
Ishikawa--Kitaoka~\cite{IshikawaKitaoka1998},
Chen--Kitaoka--Yu~\cite{ChenKitaokaYu2000},
Roskam~\cite{Roskam2000,Roskam2002}, and
Kataoka~\cite{Kataoka2001,Kataoka2003}.  Near-primitive-root and
character-sum formulations of primitive-root densities are treated in
\cite{LenstraMoreeStevenhagen2014,Palenstijn2014}.  In the present
setting the relevant point is the norm-one element
\[
  \eta=-\varphi^{2}
\]
on the torus attached to \(\mathbb Q(\sqrt5)\).  The full-rank case
\(i_p=1\) is governed, under GRH, by Chen's density theorem for
primitive points on one-dimensional tori~\cite{Chen2002}, while the
inert index-one/index-two union is identified in Lemma~\ref{lem:2adic}
with Roskam's inert maximal-order set for the fundamental unit
\(\varphi\).  The present paper uses this density framework only for
contextual information; the proofs of Theorems~\ref{thm:A} and
\ref{thm:B} are independent of GRH.
\end{remark}

\begin{remark}[Relation to the author's series]
The present paper is part of a sequence of works on rank-block
decompositions and character sums in Fibonacci and Lucas sequences.
The paper on half-period involutions~\cite{Ghandali2026HalfPeriod} treats exact
cancellation mechanisms in inert Lucas sequences, and the full-rank
paper~\cite{Ghandali2026Preprint} evaluates \(S(p)\) at \(i_p=1\).  The present paper
contributes the rank-index-two reductions
\[
  S(p)=\frac12\Jac(\cdot)
\]
together with the resulting explicit evaluations in
Theorems~\ref{thm:A} and~\ref{thm:B}.  Rank-multiplier and reflection
identities used here are classical in origin, going back to
Lucas~\cite{Lucas1878} and later treatments of Lucas sequences such as
Somer~\cite{Somer1980}; see also~\cite{Ghandali2026HalfPeriod}.  No proof in the present
paper depends on an unpublished manuscript.
\end{remark}

% =====================================================================
\section{Further problems and directions}\label{sec:open}
% =====================================================================

\begin{remark}[Signed inert two-square law]
The sign in the inert case \(i_p=2\), \(p\equiv1\pmod8\) is not an
open refinement: Theorem~\ref{thm:B} determines it by the congruence
\(S(p)\equiv-a^{*}\pmod5\).  The companion verifier (\texttt{ghandali\_verify\_index2.py}, version~1.0.1) checks this signed
law directly (zero mismatches on all \(14{,}591\) eligible primes below
\(2\cdot10^{6}\)).
\end{remark}

\begin{problem}[Index four]\label{prob:index-four}
For split primes with \(i_p=4\), the branch \(p\equiv5\pmod8\) is
explained by the reflection vanishing mechanism: in the computation up
to \(2\cdot10^{6}\), the identity
\[
  S(p)=0
\]
held with zero mismatches on \(6697\) such primes.

The remaining split branch \(p\equiv1\pmod8\) appears to contain new
structure.  In the same range, \(S(p)\) was odd for all \(3359\) tested
primes, and the maximum observed value of
\[
  \frac{|S(p)|}{\sqrt p}
\]
was
\[
  0.502683773697118.
\]
A natural next step is to seek an analogue of
Proposition~\ref{prop:split-red} involving quartic Jacobsthal sums or
Jacobi sums.
\end{problem}

\begin{problem}[Index three]\label{prob:index-three}
Preliminary computations suggest that split primes of rank index three
may be governed by cubic character sums. Establishing a precise analogue
of the index-two Jacobsthal reduction remains open.
\end{problem}


\begin{problem}[Other Lucas sequences]\label{prob:other-lucas}
The reductions in Propositions~\ref{prop:split-red} and
\ref{prop:inert-red} use the identities
\[
  \varphi\psi=-1
\]
and the associated Cayley parametrisation.  They therefore extend
formally to Lucas sequences \(U_n(P,-1)\), for which the discriminant
\(5\) is replaced by \(P^2+4\).  The final closed evaluations in
Theorem~\ref{thm:A}, however, rely on the special quartic character
identity for the Fibonacci unit.  A complete analogue for general \(P\)
would require a corresponding evaluation of the quartic character of
the relevant fundamental unit.
\end{problem}

\begin{problem}[Arbitrary fixed index]\label{prob:arbitrary-index}
For a general fixed index \(t\ge1\), can the character sum
\[
  S(p)=\sum_{j=1}^{z(p)}\chi(F_j)
\]
be expressed systematically in terms of Jacobi sums or Jacobsthal sums
of order \(t\) or \(2t\)?  The full-rank case \(t=1\) and the present
rank-index-two case suggest that fixed-index Fibonacci character sums
should admit a uniform description through finite-field character sums.
\end{problem}

% =====================================================================
\section*{Declarations}
% =====================================================================
\noindent
No external funding was received for this work.  The author declares no
competing interests.  Computational verification was performed with the
companion verifier; the versioned source code, machine-readable output
files, and run log are available in the reproducibility
archive~\cite{ghandaliIndex2Archive}.

\bibliographystyle{abbrvnat}
\bibliography{ghandali_index2}

\end{document}
